% ============================================================
%  TERM PROJECT – PROPOSAL (template)
%  Databases · THGA Bochum
%
%  How to use:
%    1. Copy this folder into your own repository.
%    2. Replace every  <<< ... >>>  placeholder with your content.
%    3. Build with  `make`  (see the Makefile / README).
%    4. Send the resulting PDF to the lecturer and wait for approval.
% ============================================================
\documentclass[11pt,a4paper]{article}

\usepackage{thga-db}
\usepackage{caption}   % enables \captionof inside minipages
\usepackage{float}
\usepackage{placeins}

% ----- Placeholder figure (delete once you insert a real image) ----
% Draws a labelled box you can swap for \includegraphics later.
% Usage: \placeholder{width}{height}{label text}
\usetikzlibrary{positioning}
\newcommand{\placeholder}[3]{%
  \begin{tikzpicture}
    \node[draw=THGABlau!45, fill=THGABlau!5, very thick, rounded corners=2pt,
          minimum width=#1, minimum height=#2, align=center, inner sep=6pt,
          text=THGABlau!75, font=\small\itshape] {#3};
  \end{tikzpicture}%
}

% ----- Fill in your metadata -------------------------------
\renewcommand{\thgaKurs}{Databases}
\renewcommand{\thgaDatum}{Summer Term 2026}
\newcommand{\authorName}{Daniel Brendebach}
\newcommand{\projectTitle}{Filament Management Software}

\begin{document}
\thispagestyle{empty}
\thgaTitel{Term Project – Proposal}{\projectTitle}

\begin{infobox}
\textbf{Author:} \authorName \\
\textbf{Module:} Introduction to Database Management Systems · THGA Bochum \\
\textbf{Submitted to:} Stephan Bökelmann · \href{mailto:sboekelmann@ep1.rub.de}{sboekelmann@ep1.rub.de} \\
\textbf{Target size:} about 6 pages
\end{infobox}


% ============================================================
\section{Application domain and motivation}
% ============================================================

I would like to develop an application that helps me manage the filament inventory for my 3D printer.
Over time, I have collected many filament spools from different manufacturers, in various materials and colors.
As the number of spools has grown, it has become increasingly difficult to keep track of which filament is available, how much material is left on each spool, and which spool should be used for a particular print.

The primary user of this application will be myself, but the software should also be designed in a way that allows other hobbyists or makers with similar requirements to use it.

The goal is to replace manual notes or spreadsheets with a structured database that stores all relevant information about the filament inventory.
The application should allow the user to register different filament types, including the manufacturer, material, color, and other product information.
For every physical spool, the current remaining weight should be stored, enabling an accurate overview of the available stock.

In addition, the application should maintain a collection of 3D print models, including the estimated amount of filament required and the material needed for each model.
One of the application's key features will be an intelligent spool selection process.
Whenever a model is selected for printing, the application should automatically determine the most suitable spool by choosing the spool with the smallest remaining amount of filament that is still sufficient to complete the print.
This minimizes filament waste and helps consume partially used spools before opening new ones.
To keep the inventory accurate over time, every completed print will be recorded in a print history.

The application will automatically reduce the remaining weight of the selected spool by the amount of filament used and store information about the printed model, the spool used, the print date, and the actual filament consumption.
This creates a complete history of all print jobs and makes it possible to analyze filament usage over time.

Overall, the project demonstrates the design and implementation of a relational database application that combines inventory management with automated decision-making while maintaining data consistency and referential integrity.
% ============================================================
\section{Data model -- ER sketch}
% ============================================================

\begin{figure}[ht]
  \centering
  \includegraphics[width=0.85\linewidth]{er-sketch.png}
  \caption{ER sketch of the filament management application data model: entity types, key
    attributes, relationships and cardinalities.}
  \label{fig:er}
\end{figure}

Figure~\ref{fig:er} shows the Entity-Relationship model of the filament management application.
The central entity types are FilamentType, FilamentRoll, PrintModel, and PrintJob.
FilamentType represents the general description of a filament product, including its manufacturer, material, and color.
Each filament type can have multiple physical spools (FilamentRoll), while each spool belongs to exactly one filament type. 
The entity FilamentRoll stores inventory-related information such as the remaining weight and purchase date.
The entity PrintModel represents a printable object and stores general information such as the model name, estimated filament consumption, and estimated printing time. A PrintJob represents an actual printing process of a specific model and stores information such as the printing date and whether the print was successful.
The relationship between PrintJob and FilamentRoll is a many-to-many (N:M) relationship. A single print job can consume multiple filament rolls, for example when a print uses multiple colors or when one spool is not sufficient to complete the print. At the same time, one filament roll can be used in multiple print jobs until the material is depleted.

% ============================================================
\section{From model to schema}
% ============================================================

\begin{itemize}
  \item \texttt{manufacturer}(\underline{manufacturer\_id}, name)
  
  \item \texttt{material}(\underline{material\_id}, name)
  
  \item \texttt{filament\_type}(\underline{filament\_type\_id}, 
  \emph{manufacturer\_id}, \emph{material\_id}, product\_name, color, diameter)
  
  \item \texttt{filament\_roll}(\underline{roll\_id}, 
  \emph{filament\_type\_id}, remaining\_weight, purchase\_date, location)
  
  \item \texttt{print\_model}(\underline{model\_id}, 
  name, estimated\_filament, estimated\_time, notes)
  
  \item \texttt{print\_job}(\underline{job\_id}, 
  \emph{model\_id}, print\_date, successful, notes)
  
  \item \texttt{print\_job\_roll}(\underline{job\_id}, 
  \underline{\emph{roll\_id}}, filament\_used)
\end{itemize}

\textbf{Normalisation}

The schema is normalized to at least Third Normal Form (3NF), because each
relation represents a single entity or relationship and every non-key attribute
depends only on the primary key of its relation. Redundant information such as
manufacturer, material, and filament properties is stored only once and
referenced using foreign keys.

\textbf{N to M}

The many-to-many relationship between \texttt{print\_job} and
\texttt{filament\_roll} is resolved by the associative relation
\texttt{print\_job\_roll}, which stores the relationship-specific attribute
\texttt{filament\_used}. This avoids data duplication and maintains referential
integrity.

\clearpage
% ============================================================
\section{API design}
% ============================================================

<<< List the planned FastAPI endpoints. Mark where a JOIN / an aggregation
sits, and which endpoints require the X-API-Key. >>>

\renewcommand{\arraystretch}{1.2}
\begin{tabularx}{\linewidth}{@{}p{2cm}p{4cm}X@{}}
\toprule
\textbf{Method} & \textbf{Path} & \textbf{Purpose} \\
\midrule

\texttt{GET} &
\texttt{/filament/types} &
Reads all filament types including manufacturer and material information
(JOIN between \texttt{filament\_type}, \texttt{manufacturer} and
\texttt{material}). \\

\texttt{GET} &
\texttt{/filament/rolls} &
Reads all filament rolls including their filament type information
(JOIN between \texttt{filament\_roll} and \texttt{filament\_type}). \\

\texttt{GET} &
\texttt{/filament/stock} &
Reads the current filament inventory grouped by filament type and calculates
the available amount using aggregation (\texttt{SUM()} and \texttt{GROUP BY}). \\

\texttt{POST} &
\texttt{/filament/rolls} &
Creates a new filament roll entry with manufacturer, material and weight
information (requires \texttt{X-API-Key}). \\

\texttt{GET} &
\texttt{/models} &
Reads all available 3D print models stored in the database. \\

\texttt{POST} &
\texttt{/models} &
Creates a new 3D print model entry
(requires \texttt{X-API-Key}). \\

\texttt{GET} &
\texttt{/print/jobs} &
Reads the print history including model information and consumed filament
rolls (JOIN between \texttt{print\_job}, \texttt{print\_model},
\texttt{print\_job\_roll} and \texttt{filament\_roll}). \\

\texttt{POST} &
\texttt{/print/jobs} &
Creates a new print job, selects suitable filament rolls, updates the remaining
filament weight and stores the consumption history
(requires \texttt{X-API-Key}). \\


\bottomrule
\end{tabularx}

\clearpage

% ============================================================
\section{Frontend prototype (hand sketches)}
% ============================================================

The frontend will be a desktop application built using Tkinter and Python.

% Each \includegraphics below still points at the dummy example-figure.jpg.
% Swap in your own file per screen, e.g. {fe-connect.jpg}, and keep the \captionof.
\begin{figure}[ht]
\centering
\begin{minipage}[t]{0.47\linewidth}\centering
  \includegraphics[width=\linewidth]{Connection.png}
  \captionof{figure}{Startup connection dialog: API URL + X-API-Key.}
  \label{fig:fe-connect}
\end{minipage}\hfill
\begin{minipage}[t]{0.47\linewidth}\centering
  \includegraphics[width=\linewidth]{filament_rolls.png}
  \captionof{figure}{Main screen: filament roll table, calls \texttt{GET /filament/rolls}.}
  \label{fig:fe-main}
\end{minipage}

\vspace{1em}

\begin{minipage}[t]{0.47\linewidth}\centering
  \includegraphics[width=\linewidth]{print_models.png}
  \captionof{figure}{Detail form: print models table, calls \texttt{GET /models}.}
  \label{fig:fe-detail}
\end{minipage}\hfill
\begin{minipage}[t]{0.47\linewidth}\centering
  \includegraphics[width=\linewidth]{print_history.png}
  \captionof{figure}{Detail form: print history table, calls \texttt{GET /print/jobs}.}
  \label{fig:fe-second}
\end{minipage}
\end{figure}

\FloatBarrier

Figures~\ref{fig:fe-connect}--\ref{fig:fe-second} show the four screens of the
prototype. Figure~\ref{fig:fe-connect} shows the connection dialog that
collects the API URL and X-API-Key at startup ; Figure~\ref{fig:fe-main}
shows the main screen and the endpoint it reads from ;
Figure~\ref{fig:fe-detail} shows the print model overview and the endpoint it reads from ;
 and Figure~\ref{fig:fe-second} shows the print history and the endpoint it reads from.



% ============================================================
\section{Operation with Docker Compose}
% ============================================================

The application is deployed using \texttt{docker-compose}. The compose file
starts two services:

\begin{itemize}
    \item \textbf{postgres}: A PostgreSQL database container that stores all
    application data.
    \item \textbf{api}: A FastAPI backend container that provides the REST API
    and communicates with the PostgreSQL database.
\end{itemize}

Database credentials and other configuration parameters are stored in a
\texttt{.env} file. This includes the PostgreSQL username, password, database
name, host, port and the API authentication key. Using a separate
\texttt{.env} file keeps sensitive information out of the source code and
allows different configurations for development and deployment.

On the lecture server, the application is deployed using Docker Compose. After
cloning the Git repository, the containers are started with
\texttt{docker compose up -d}. The FastAPI backend is then available via its
configured port, while the PostgreSQL database is only accessible from within
the Docker network. Persistent database storage is provided through a Docker
volume so that the stored data remains available even if the containers are
recreated.

% ============================================================
\section{Scope, milestones and risks}
% ============================================================

\textbf{Scope.}
The application consists of \textbf{7 database tables} (\texttt{manufacturer},
\texttt{material}, \texttt{filament\_type}, \texttt{filament\_roll},
\texttt{print\_model}, \texttt{print\_job}, and
\texttt{print\_job\_roll}) and \textbf{8 REST API endpoints}. The API provides
functions for managing filament, print models, print jobs, inventory
statistics, and automatic filament selection for new print jobs.

\textbf{Milestones and schedule.}

\renewcommand{\arraystretch}{1.2}
\begin{tabularx}{\linewidth}{@{}p{1.4cm}Xp{1.8cm}@{}}
\toprule
\textbf{Week} & \textbf{Milestone} & \textbf{Est. hours} \\
\midrule
1 &
Database design completed (ER model, relational schema, PostgreSQL database,
Docker Compose, seed data) &
8 \\

2 &
Implementation of the FastAPI backend, database connection and CRUD endpoints
for filament types, rolls and print models &
10 \\

3 &
Implementation of print jobs, automatic filament selection, inventory updates
and API authentication using \texttt{X-API-Key} &
10 \\

4 &
Frontend development including inventory management, print history and filament
selection &
8 \\

5 &
Testing, bug fixing, documentation, packaging as a Debian package and final
project presentation &
4 \\

\bottomrule
\end{tabularx}

\textbf{Risks.}
The automatic filament selection algorithm may become more complex than
expected, especially when multiple filament rolls are required for a single
print job. Another risk is ensuring data consistency when updating the remaining
filament weight after completed print jobs, which requires careful transaction
handling to avoid inconsistent inventory data.

% ============================================================
\section{Acceptance criteria and testing}
% ============================================================

The system is considered complete when all planned database functions, API
endpoints and business rules are implemented and tested successfully. Every
acceptance criterion can be verified by automated or manual tests.

\textbf{Acceptance criteria.}
\begin{itemize}[topsep=2pt]
    \item The database enforces all defined constraints (PRIMARY KEY, FOREIGN
    KEY, NOT NULL and CHECK). Invalid records are rejected.
    
    \item The API correctly creates, updates and retrieves filament rolls,
    filament types, print models and print jobs.

    \item The endpoint \texttt{GET /filament/stock} returns the correct
    aggregated remaining filament weight for every filament type (aggregation).

    \item The endpoint \texttt{GET /filament/types} correctly returns filament
    types together with their manufacturer and material information (JOIN).

    \item When a new print job is created, the system automatically selects the
    smallest suitable filament roll (or multiple rolls if necessary) and
    correctly updates the remaining filament weight.

    \item All write endpoints reject requests without a valid
    \texttt{X-API-Key} by returning HTTP status code 401 or 403.
\end{itemize}

\textbf{Test plan.}
\begin{itemize}[topsep=2pt]

    \item \textbf{Database:}
    Verify PRIMARY KEY, FOREIGN KEY, NOT NULL and CHECK constraints by trying
    to insert invalid records and confirming that PostgreSQL rejects them.

    \item \textbf{API:}
    Test all REST endpoints and look for error
    cases such as invalid input, missing resources and unauthorized access.

    \item \textbf{Business rule:}
    Verify that the filament selection algorithm always chooses the smallest
    filament roll with sufficient remaining weight. Additionally, test the case
    where no single roll contains enough filament, requiring the system to use
    multiple rolls for one print job while correctly updating the remaining
    weight of each roll.

\end{itemize}



\vspace{1em}
\begin{infobox}
\textbf{Effort budget:} the whole term project (system, documentation and video)
should fit into roughly \textbf{40\,h of work} over at most \textbf{2 months}.
Keep your scope within that budget.
\end{infobox}

\end{document}
